\documentclass[11pt]{article}
\usepackage[margin=2.5cm]{geometry}
\usepackage{algorithm}
\usepackage{algpseudocode}
\usepackage{amsmath}
\usepackage{amssymb}
\usepackage{hyperref}

\title{Hivatalos válaszok Siket István tanár úr részére}
\author{}
\date{}

\begin{document}
\maketitle

\section{1. kérdés}

\paragraph {Kérdés}
A “Fordítási Egységek Topologikus Rendezése” alfejezetben a TU-kat topologikusan rendezik. Hogyan kezeli azt az esetet, ha kör van a hívási gráfban, és nincs topologikus rendezés?

\paragraph {Válasz}

Az eszköz fordítási egységek közötti analízise során Kahn algoritmusát \cite{Kahn1962} használja, hogy egyszerre detektálja a függőségi gráfban a körök megjelenését, illetve hogy azon részgráfban amely körmentes gráfot alkot, megadjon egy topologikus rendezést.

These dependencies form a directed graph $G = (V, E)$ where $V$ is the set of TUs and $(A, B) \in E$ means $A$ depends on $B$.

This graph may contain cycles (back-edges). For example, TU~$A$ calls a function in TU~$B$, and TU~$B$ calls a function in TU~$A$. A pure topological sort is undefined for cyclic graphs, and analyzing TUs in any linear order will leave some dependencies unresolved after a single pass.

\section{Solution: Three-Level Approach}

The tool handles cycles at three levels:

\subsection{Level 1: Best-Effort Topological Sort (Kahn's Algorithm)}

The tool uses Kahn's algorithm to produce a processing order that respects as many dependencies as possible. Nodes involved in cycles are appended at the end in arbitrary order.

\begin{algorithm}[H]
\caption{Best-Effort Topological Sort}
\begin{algorithmic}[1]
\Procedure{TopologicalSort}{$G = (V, E)$}
    \State $\text{InDegree}[v] \gets |\{u : (u, v) \in E\}|$ for all $v \in V$
    \State $Q \gets \{v \in V : \text{InDegree}[v] = 0\}$
    \State $\text{Result} \gets []$
    \While{$Q \neq \emptyset$}
        \State $v \gets Q.\text{pop}()$
        \State $\text{Result}.\text{append}(v)$
        \ForAll{$(v, w) \in E$}
            \State $\text{InDegree}[w] \gets \text{InDegree}[w] - 1$
            \If{$\text{InDegree}[w] = 0$}
                \State $Q.\text{add}(w)$
            \EndIf
        \EndFor
    \EndWhile
    \State $\text{HasCycle} \gets \exists\, v \in V : \text{InDegree}[v] > 0$
    \ForAll{$v \in V$ where $\text{InDegree}[v] > 0$} \Comment{Cycle nodes}
        \State $\text{Result}.\text{append}(v)$
    \EndFor
    \State \Return $(\text{Result},\; \neg\text{HasCycle})$
\EndProcedure
\end{algorithmic}
\end{algorithm}

The result is reversed before use (dependencies before dependents). Cycle nodes are analyzed with potentially incomplete information (defaulting to \texttt{Unknown})---but this is resolved by Level~2.

\subsection{Level 2: Fixed-Point Iteration}

The analysis is run repeatedly on a \textbf{worklist of Translation Units (TUs)} until a fixed point is reached. This resolves circular dependencies across TUs.

\begin{algorithm}[H]
\caption{Fixed-Point Analysis}
\begin{algorithmic}[1]
\Procedure{RunUntilFixedPoint}{$\text{Phase}, \text{TUs}$}
    \State $\text{WorkList} \gets \text{TUs}$ \Comment{Initially all TUs}
    \While{$\text{WorkList} \neq \emptyset$}
        \State $\text{Changed} \gets \textbf{false}$
        \ForAll{$\text{TU} \in \text{WorkList}$}
            \State Run $\text{Phase}$ on $\text{TU}$, updating shared \texttt{USRToExceptionMap}
            \If{shared state was modified}
                \State $\text{Changed} \gets \textbf{true}$
            \EndIf
        \EndFor
        \If{$\neg\text{Changed}$}
            \State \textbf{break}
        \EndIf
        \State $\text{WorkList} \gets \text{dependents of changed TUs}$
    \EndWhile
\EndProcedure
\end{algorithmic}
\end{algorithm}

\paragraph{Lattice and Convergence.} The exception state lattice is finite and monotonic:
\[
\text{NotThrowing} \sqsubseteq \text{Unknown} \sqsubseteq \text{Throwing}
\]
Each iteration only moves a function's state upward (from \texttt{NotThrowing} toward \texttt{Throwing}). Fixed point is guaranteed as states never revert.

\subsection{Level 3: Cycle Detection and Reporting}

After analysis, the tool explicitly detects TU cycles using a Depth-First Search (DFS) with a recursion stack to report them as diagnostics.

\begin{algorithm}[H]
\caption{TU Cycle Detection (DFS)}
\begin{algorithmic}[1]
\Procedure{DetectCycles}{$G = (V, E)$}
    \State $\text{Visited} \gets \emptyset$, $\text{RecStack} \gets \emptyset$, $\text{Cycles} \gets []$
    \ForAll{$v \in V$ not yet visited}
        \State \Call{DFS}{$v, \text{Visited}, \text{RecStack}, \text{Path}$}
    \EndFor
    \State \Return $\text{Cycles}$
\EndProcedure
\end{algorithmic}
\end{algorithm}

\begin{thebibliography}{9}
\bibitem{Kahn1962}
Arthur B. Kahn. 1962. Topological sorting of large networks. \textit{Commun. ACM} 5, 11 (Nov. 1962), 558–562. \url{https://doi.org/10.1145/368996.369025}
\end{thebibliography}

Cyclic SCCs are reported to the user as diagnostics and written to a \texttt{.dot} file for visualization. This is purely informational---the analysis has already handled cycles via fixed-point iteration.

\section{Summary}

\begin{center}
\begin{tabular}{lll}
\hline
\textbf{Level} & \textbf{Mechanism} & \textbf{Purpose} \\
\hline
1. Topological sort & Kahn's algorithm & Best-effort processing order \\
2. Fixed-point iteration & Re-analyze until stable & Resolve circular dependencies \\
3. Cycle detection & Tarjan's SCC algorithm & Report cycles to user \\
\hline
\end{tabular}
\end{center}

Back-edges (cycles) do not break correctness. They increase the number of fixed-point iterations needed (typically from 1 to 2--3) but the analysis is guaranteed to converge due to the monotonic exception state lattice.

\end{document}
